% Chapter 3: Automata, Languages, and Computability
% Rewritten against the design-lab Lean 4 kernel (FLT_Proofs.Computation, FLT_Proofs.Complexity.MindChange):
% every model typed against the verified computability substrate; full proofs for the
% pumping separation, the subset construction, and the limit-lemma direction.
%   dfa, nfa, pushdown_automaton, turing_machine, cfg = Profile A (vocabulary)
%   regular_language, cfl, re_set, partial_computable = Profile A (one-line)
%   godel_numbering = Profile A-plus (brief but load-bearing for Gold)
%   restricts/extends_grammar edge pair = Profile C (Revelation: conceptual centerpiece)
%   limiting_recursion + delta_02_class = Profile G (Framework Bridge to Gold)

\chapter{Automata, Languages, and Computability}
\label{ch:automata}

This chapter assembles the automata-theoretic and computability-theoretic vocabulary required by Gold's model of identification in the limit (\Cref{ch:gold}). Much of the material is standard; the reader who needs the full classical development should consult Hopcroft and Ullman~\cite{HopcroftUllman1979} or Sipser~\cite{Sipser2013}. Every model defined here is typed in the companion Lean~4 development's computability substrate, and three results---the pumping separation of regular from context-free languages, the subset construction, and one direction of the limit lemma---are proved in full.

Two things here are not standard vocabulary. First, the distinction between \emph{restricts} and \emph{extends\_grammar} edges---two ways one computational model can generalize another, one conservative and one genuinely new---is illustrated through the NFA/DFA and PDA/DFA pairs (\S\ref{sec:restricts-extends}); the distinction pervades the entire knowledge graph. Second, limiting recursion connects computability to learning by placing Gold's identification in the limit at a precise level of the arithmetical hierarchy (\S\ref{sec:bridge}).

% ================================================================
% SECTION 1: The Chomsky Hierarchy
% ================================================================

\section{The Chomsky Hierarchy}
\label{sec:chomsky}

The Chomsky hierarchy classifies formal languages by the computational resources needed to recognize them. Figure~\ref{fig:chomsky} displays the containment structure; the definitions follow. A \emph{formal language} over an alphabet $\Sigma$ is a set of finite words, $L \subseteq \Sigma^\ast$.\formal{FLT_Proofs.Computation.FormalLanguage}

\begin{figure}[ht]
\centering
\begin{tikzpicture}[
    every node/.style={font=\small},
    level/.style={draw, thick, ellipse, minimum height=1.2cm, align=center}
]
% Nested ovals, outermost first
\node[level, minimum width=10cm, minimum height=5.8cm, fill=layergray!60] (re) at (0,0) {};
\node[font=\scriptsize, anchor=north] at (re.north) {Recursively enumerable (Type 0)};

\node[level, minimum width=7.8cm, minimum height=4.2cm, fill=defblue!6] (rec) at (0,-0.3) {};
\node[font=\scriptsize, anchor=north] at (rec.north) {Recursive (decidable)};

\node[level, minimum width=5.6cm, minimum height=2.8cm, fill=edgeorange!8] (cfl) at (0,-0.5) {};
\node[font=\scriptsize, anchor=north] at (cfl.north) {Context-free (Type 2)};

\node[level, minimum width=3.2cm, minimum height=1.4cm, fill=sepgreen!10] (reg) at (0,-0.6) {};
\node[font=\scriptsize] at (reg.center) {Regular (Type 3)};

% Recognizer labels on the right
\node[font=\scriptsize\itshape, notegray, anchor=west] at (4.2,2.2) {Turing machine};
\node[font=\scriptsize\itshape, notegray, anchor=west] at (3.2,0.8) {TM (halting)};
\node[font=\scriptsize\itshape, notegray, anchor=west] at (2.2,-0.1) {Pushdown automaton};
\node[font=\scriptsize\itshape, notegray, anchor=west] at (1.2,-0.6) {DFA / NFA};
\end{tikzpicture}
\caption{The Chomsky hierarchy. Each level is strictly contained in the next. The recognizer type is indicated at right. All containments are proper.}
\label{fig:chomsky}
\end{figure}

\subsection{Recognizers}

\begin{defn}[Deterministic Finite Automaton]
\label{def:dfa}
A \emph{DFA} is a 5-tuple $M = (Q, \Sigma, \delta, q_0, F)$ where $Q$ is a finite state set, $\Sigma$ a finite input alphabet, $\delta : Q \times \Sigma \to Q$ a total transition function, $q_0 \in Q$ the start state, and $F \subseteq Q$ the accepting states. $M$ accepts $w \in \Sigma^\ast$ if the unique run on $w$ ends in $F$; the accepted language is $\{w : \hat\delta(q_0, w) \in F\}$, where $\hat\delta$ folds $\delta$ along $w$.\formal{FLT_Proofs.Computation.DFA'}
\end{defn}

\begin{defn}[Nondeterministic Finite Automaton]
\label{def:nfa}
An \emph{NFA} is a 5-tuple identical to a DFA except that $\delta : Q \times (\Sigma \cup \{\varepsilon\}) \to \mathcal{P}(Q)$ maps each state--symbol pair to a \emph{set} of successors and $\varepsilon$-transitions are permitted; $M$ accepts $w$ if \emph{some} run ends in $F$.\formal{FLT_Proofs.Computation.NFA'}
\end{defn}

\begin{defn}[Pushdown Automaton]
\label{def:pda}
A \emph{pushdown automaton} (PDA) is a 7-tuple $M = (Q, \Sigma, \Gamma, \delta, q_0, Z_0, F)$ where $\Gamma$ is a finite stack alphabet, $Z_0 \in \Gamma$ the initial stack symbol, and $\delta : Q \times (\Sigma \cup \{\varepsilon\}) \times \Gamma \to \mathcal{P}(Q \times \Gamma^\ast)$ maps each state--input--stack-top triple to a set of (state, stack-replacement) pairs. The PDA has an unbounded LIFO stack and accepts by final state or empty stack.\formal{FLT_Proofs.Computation.PDA}
\end{defn}

\begin{defn}[Turing Machine]
\label{def:tm}
A \emph{Turing machine} is a 7-tuple $M = (Q, \Sigma, \Gamma, \delta, q_0, q_{\mathrm{acc}}, q_{\mathrm{rej}})$ where $\Gamma \supseteq \Sigma$ is the tape alphabet (with a blank), and $\delta : Q \times \Gamma \to Q \times \Gamma \times \{L, R\}$ is a partial transition function specifying the next state, the symbol to write, and head movement; $M$ accepts $w$ if it halts in $q_{\mathrm{acc}}$.\formal{FLT_Proofs.Computation.TuringMachine}
\end{defn}

\begin{defn}[Context-Free Grammar]
\label{def:cfg}
A \emph{context-free grammar} (CFG) is a 4-tuple $G = (V, \Sigma, R, S)$ with $V$ a finite variable set, $\Sigma$ the terminal alphabet, $R \subseteq V \times (V \cup \Sigma)^\ast$ a finite rule set, and $S \in V$ the start variable; $G$ generates $L(G) = \{w \in \Sigma^\ast : S \Rightarrow^\ast w\}$, the reflexive--transitive closure of one-step derivation.\formal{FLT_Proofs.Computation.CFG}
\end{defn}

\subsection{Language classes}

A language $L \subseteq \Sigma^\ast$ is \emph{regular} if accepted by some DFA;\formal{FLT_Proofs.Computation.IsRegularLang} \emph{context-free} if generated by some CFG (equivalently, accepted by some PDA);\formal{FLT_Proofs.Computation.IsContextFree} \emph{recursively enumerable} (r.e.) if accepted by some Turing machine;\formal{FLT_Proofs.Computation.RESet} and \emph{recursive} (decidable) if accepted by a Turing machine that halts on every input. A function $f : \N \rightharpoonup \N$ is \emph{partial computable} if some Turing machine outputs $f(n)$ when $f(n)$ is defined and diverges otherwise.\formal{FLT_Proofs.Computation.PartialComputable}

\subsection{Separating the levels by pumping}

The containments of \cref{fig:chomsky} are proper, and the regular/context-free gap drives the edge distinction of \S\ref{sec:restricts-extends}. We prove it from a single combinatorial lemma, whose only input is that a DFA has finitely many states.

\begin{lemma}[Pumping lemma for regular languages {\cite{HopcroftUllman1979}}]
\label{lemma:pumping}
For every regular $L$ there is a length $p$ such that every $w \in L$ with $|w| \geq p$ factors as $w = xyz$ with $|xy| \leq p$, $|y| \geq 1$, and $xy^iz \in L$ for all $i \geq 0$.
\end{lemma}

\begin{proof}
Let $M$ be a DFA accepting $L$, with $p = |Q|$ states. Take $w \in L$ with $|w| \geq p$. The run of $M$ on $w$ passes through the states $r_0, r_1, \ldots, r_{|w|}$ with $r_0 = q_0$; among the first $p+1$ of these, two coincide by the pigeonhole principle, say $r_j = r_k$ with $j < k \leq p$. Set $x = w_{1\ldots j}$, $y = w_{j+1\ldots k}$, $z = w_{k+1\ldots |w|}$. Then $|xy| = k \leq p$ and $|y| = k - j \geq 1$. Reading $y$ takes state $r_j$ back to itself, so reading $y$ any number of times leaves the run in $r_j$; hence $xy^iz$ drives $M$ from $q_0$ to the same accepting state $r_{|w|}$ for every $i \geq 0$, and $xy^iz \in L$.
\end{proof}

\begin{prop}[The language $\{a^n b^n\}$ is context-free but not regular]
\label{prop:anbn}
$L = \{a^n b^n : n \geq 0\}$ is context-free and not regular; hence $\mathrm{REG} \subsetneq \mathrm{CFL}$.
\end{prop}

\begin{proof}
$L$ is generated by the CFG $S \to aSb \mid \varepsilon$ (equivalently, accepted by a PDA that pushes on each $a$ and pops on each $b$), so $L$ is context-free. Suppose $L$ were regular, with pumping length $p$ from \cref{lemma:pumping}. Apply the lemma to $w = a^p b^p \in L$, with $|w| = 2p \geq p$. The factorization $w = xyz$ has $|xy| \leq p$, so $xy$ lies within the initial block of $a$'s and $y = a^k$ for some $k \geq 1$. Then $xy^2z = a^{p+k} b^p$, which has more $a$'s than $b$'s and so is not in $L$, contradicting the lemma. Therefore $L$ is not regular, and since $L \in \mathrm{CFL}$, the containment $\mathrm{REG} \subseteq \mathrm{CFL}$ is proper.
\end{proof}

\subsection{Closure properties}

The classes differ sharply in closure. Table~\ref{tab:closure} records the facts used in later chapters; the proofs are standard~\cite{HopcroftUllman1979}, and we give the two that recur. Regular languages are closed under union by the product construction; context-free languages are \emph{not} closed under intersection.

\begin{prop}[Two closure facts]
\label{prop:closure}
(i) If $L_1, L_2$ are regular then so is $L_1 \cup L_2$. (ii) Context-free languages are not closed under intersection.
\end{prop}

\begin{proof}
(i) Let $M_1 = (Q_1,\Sigma,\delta_1,s_1,F_1)$ and $M_2 = (Q_2,\Sigma,\delta_2,s_2,F_2)$ accept $L_1, L_2$. The product DFA on state set $Q_1 \times Q_2$ with transition $\delta((q_1,q_2),a) = (\delta_1(q_1,a),\delta_2(q_2,a))$, start $(s_1,s_2)$, and accepting set $\{(q_1,q_2): q_1 \in F_1 \text{ or } q_2 \in F_2\}$ simulates both machines in parallel and accepts exactly $L_1 \cup L_2$.

(ii) The languages $L_1 = \{a^n b^n c^m : n,m \geq 0\}$ and $L_2 = \{a^m b^n c^n : n,m \geq 0\}$ are each context-free (a single matched block, the other free), but $L_1 \cap L_2 = \{a^n b^n c^n : n \geq 0\}$, which is not context-free---the pumping lemma for context-free languages forces any pumped factor to disturb the three-way balance. Hence the context-free languages are not closed under intersection.
\end{proof}

\begin{table}[ht]
\centering
\small
\caption{Closure properties of the Chomsky hierarchy classes. $\checkmark$ = closed, $\times$ = not closed. Union for regular and intersection for context-free are proved in \cref{prop:closure}; the remaining entries are standard~\cite{HopcroftUllman1979}.}
\label{tab:closure}
\begin{tabular}{lccccc}
\toprule
\textbf{Class} & \textbf{Union} & \textbf{Intersection} & \textbf{Complement} & \textbf{Concat.} & \textbf{Kleene $*$} \\
\midrule
Regular       & $\checkmark$ & $\checkmark$ & $\checkmark$ & $\checkmark$ & $\checkmark$ \\
Context-free  & $\checkmark$ & $\times$     & $\times$     & $\checkmark$ & $\checkmark$ \\
Recursive     & $\checkmark$ & $\checkmark$ & $\checkmark$ & $\checkmark$ & $\checkmark$ \\
R.E.          & $\checkmark$ & $\checkmark$ & $\times$     & $\checkmark$ & $\checkmark$ \\
\bottomrule
\end{tabular}
\end{table}

% ================================================================
% SECTION 2: restricts vs extends_grammar (Profile C — Revelation)
% ================================================================

\section{Two Kinds of Generalization: Restricts vs.\ Extends}
\label{sec:restricts-extends}

The Chomsky hierarchy is a containment chain, but its edges are not all of one type. Moving between levels sometimes removes a constraint (a more flexible model with no new primitive) and sometimes adds a genuinely new resource (a model with a mechanism it lacked). This distinction---conservative versus generative generalization---is a recurring structural motif of the knowledge graph, and the automata hierarchy provides its first concrete instance.

\begin{defn}[Restricts]
\label{def:restricts}
Model $B$ \rel{restricts} model $A$ if $B$ is obtained from $A$ by \emph{removing a constraint}: every $A$ is already a $B$ (after trivial embedding), and the two recognize the same class of languages. The generalization is conservative---it adds representational flexibility without expanding expressive power.
\end{defn}

\begin{defn}[Extends Grammar]
\label{def:extends-grammar}
Model $B$ \rel{extends\_grammar} model $A$ if $B$ is obtained by \emph{adding a new computational primitive}: $B$ recognizes a strictly larger class. The generalization is generative---it requires a fundamentally new resource.
\end{defn}

The DFA is the anchor for both edges. Figure~\ref{fig:restricts-extends} displays the structure, and the two propositions below supply the rigorous content of each edge.

\begin{figure}[ht]
\centering
\begin{tikzpicture}[
    box/.style={draw, thick, rounded corners, minimum width=2.4cm, minimum height=0.9cm, font=\small\bfseries, align=center},
    arr/.style={-{Stealth[length=2.5mm]}, thick},
    lbl/.style={font=\scriptsize, fill=white, inner sep=2pt, align=center}
]
\node[box, fill=sepgreen!12] (dfa) at (0,0) {DFA};
\node[box, fill=sepgreen!12] (nfa) at (5,1.5) {NFA};
\node[box, fill=edgeorange!12] (pda) at (5,-1.5) {PDA};

\draw[arr, sepgreen!70!black] (nfa) -- node[lbl, above left] {\rel{restricts}\\{\scriptsize\itshape removes determinism}\\{\scriptsize\itshape same power}} (dfa);
\draw[arr, edgeorange!80!black] (pda) -- node[lbl, below left] {\rel{extends\_grammar}\\{\scriptsize\itshape adds stack}\\{\scriptsize\itshape strictly more power}} (dfa);

% Annotations
\node[font=\scriptsize, text=sepgreen!70!black, anchor=west] at (6, 1.5) {Regular languages};
\node[font=\scriptsize, text=edgeorange!80!black, anchor=west] at (6, -1.5) {Context-free languages};
\end{tikzpicture}
\caption{Two generalizations of the DFA. The NFA removes the determinism constraint (conservative: same expressive power, \cref{thm:nfa-dfa}). The PDA adds a stack (generative: strictly more power, \cref{prop:anbn}). Both arrows point toward the DFA to indicate that the DFA is the more constrained model.}
\label{fig:restricts-extends}
\end{figure}

\paragraph{The NFA--DFA edge: conservative generalization.}
An NFA relaxes the DFA by allowing several successors per transition and $\varepsilon$-moves, removing the determinism constraint. The subset construction shows this buys no expressive power.

\begin{thm}[Subset construction {\cite{RabinScott1959}}]
\label{thm:nfa-dfa}
Every NFA with $n$ states has an equivalent DFA with at most $2^n$ states; hence $\mathrm{REG}_{\mathrm{NFA}} = \mathrm{REG}_{\mathrm{DFA}}$.
\end{thm}

\begin{proof}
Let $N = (Q, \Sigma, \Delta, S_0, F)$ be an NFA (with $S_0$ the set of start states after $\varepsilon$-closure). Define a DFA $D$ on the power set $2^Q$ by $\delta(T, a) = \bigcup_{q \in T} \Delta(q, a)$ (each set extended by $\varepsilon$-closure), start state $S_0$, and accepting sets $\{T : T \cap F \neq \varnothing\}$. By induction on $|w|$, the state $\hat\delta(S_0, w)$ of $D$ after reading $w$ is exactly the set of states $N$ can occupy after $w$: the base case is $S_0$, and the inductive step is the definition of $\delta$. Thus $D$ accepts $w$ iff some run of $N$ ends in $F$ iff $N$ accepts $w$, so $L(D) = L(N)$. The DFA has $|2^Q| = 2^n$ states.
\end{proof}

This is a \rel{restricts} edge: the NFA needs no new primitive, only a relaxed structural requirement. The cost is representational---an exponential blow-up in states---not expressive.

\paragraph{The PDA--DFA edge: generative generalization.}
A PDA adds an unbounded stack: not the removal of a constraint but the introduction of a new primitive, the ability to store and retrieve unbounded LIFO information. By \cref{prop:anbn}, the consequence is a strict increase in power---$\{a^n b^n\}$ is context-free but not regular, so $\mathrm{REG} \subsetneq \mathrm{CFL}$. This is an \rel{extends\_grammar} edge, and the gap is witnessed by the failure of the pumping lemma, a property that holds for every regular language and fails for $\{a^n b^n\}$.

\paragraph{Why this distinction matters.}
When $B$ \rel{restricts} $A$, the theory of $A$ transfers intact---same characterization theorems, same closure properties, same complexity measures. When $B$ \rel{extends\_grammar} $A$, the old theory typically breaks: new characterization theorems are needed, and closure properties change (\cref{tab:closure}: context-free languages lose closure under intersection and complement, \cref{prop:closure}). The distinction reappears in the learning setting. Relaxing a convergence requirement (behaviorally correct learning relaxing syntactic convergence) is conservative; adding a resource (a teacher who supplies counterexamples) is generative. Recognizing which is at work predicts whether existing results transfer.

% ================================================================
% SECTION 3: Computability Foundations
% ================================================================

\section{Computability Foundations for Learning}
\label{sec:computability}

Gold's model requires learners to be effective procedures that output hypotheses from an effective enumeration. This section fixes the computability vocabulary; each notion is typed in the verified development's computability substrate.

\begin{defn}[G\"odel Numbering]
\label{def:godel-numbering}
A \emph{G\"odel numbering} of a countable set is an effective bijection with $\N$.\formal{FLT_Proofs.Computation.GodelNumbering} We write $\varphi_e$ for the partial computable function with index $e$, and $W_e = \mathrm{dom}(\varphi_e)$ for the $e$-th r.e.\ set; the sequences $\varphi_0, \varphi_1, \ldots$ and $W_0, W_1, \ldots$ are universal effective enumerations of all partial computable functions and all r.e.\ sets.
\end{defn}

In Gold's framework a learner reading a text for $L$ outputs at each step a number $e_t$, interpreted through the G\"odel numbering as the conjecture $L = W_{e_t}$. Identification in the limit means the sequence $e_0, e_1, \ldots$ stabilizes on an index $e^\ast$ with $W_{e^\ast} = L$. The learner is itself a program: a computable map $M : \Sigma^\ast \to \N$ from finite data sequences to hypothesis indices, and the admissible learners are exactly the total computable such maps. This effectivity requirement separates Gold's model from a purely set-theoretic notion of convergence.

% ================================================================
% SECTION 4: The Learning-Theoretic Bridge
% ================================================================

\section{The Learning-Theoretic Bridge: Limiting Recursion and $\Delta^0_2$}
\label{sec:bridge}

The connection between computability and learning is an equivalence that identifies Gold's criterion with a precise level of the arithmetical hierarchy.

\begin{defn}[Limiting Recursion]
\label{def:limiting-recursion}
A function $f : \N \to \N$ is \emph{limiting recursive} if there is a total computable $\hat{f} : \N \times \N \to \N$ with $f(x) = \lim_{s \to \infty} \hat{f}(x, s)$ for every $x$: the approximation $\hat{f}(x,0), \hat{f}(x,1), \ldots$ stabilizes on $f(x)$, though the stabilization time is not itself computable.\formal{FLT_Proofs.Computation.LimitingRecursive}
\end{defn}

\begin{defn}[$\Delta^0_2$ Class]
\label{def:delta-02}
A set $A \subseteq \N$ is \emph{$\Delta^0_2$} if both $A$ and $\bar{A}$ are $\Sigma^0_2$, that is, of the form $\{x : \exists y\, \forall z\, R(x,y,z)\}$ for a computable $R$.\formal{FLT_Proofs.Computation.Delta02Class}
\end{defn}

The two notions coincide, and one direction is a clean unfolding of the definitions; this is the half we prove. The converse is Shoenfield's limit lemma.

\begin{prop}[Limiting recursion gives $\Delta^0_2$]
\label{prop:limit-delta}
If the characteristic function $\chi_A$ of $A \subseteq \N$ is limiting recursive, then $A$ is $\Delta^0_2$.
\end{prop}

\begin{proof}
Let $\hat{\chi} : \N \times \N \to \{0,1\}$ be a total computable approximation with $\chi_A(x) = \lim_s \hat{\chi}(x,s)$ for all $x$. Because the limit exists, the approximation is eventually constant at each $x$, so
\[
x \in A \iff \exists s_0\, \forall s \geq s_0\ \hat{\chi}(x,s) = 1, \qquad
x \notin A \iff \exists s_0\, \forall s \geq s_0\ \hat{\chi}(x,s) = 0 .
\]
The right-hand side of the first equivalence is $\Sigma^0_2$ (an $\exists\forall$ over the computable predicate $\hat{\chi}(x,s)=1$), and the second exhibits $\bar{A}$ as $\Sigma^0_2$ likewise. Both $A$ and $\bar A$ are $\Sigma^0_2$, so $A$ is $\Delta^0_2$.
\end{proof}

\begin{thm}[Shoenfield's limit lemma {\cite{Shoenfield1959}}]
\label{thm:shoenfield}
$A \subseteq \N$ is $\Delta^0_2$ if and only if $\chi_A$ is limiting recursive, equivalently iff $A$ is computable from an oracle for the halting problem.
\end{thm}

The forward direction is \cref{prop:limit-delta}; the converse extracts a computable mind-changing approximation from a $\Delta^0_2$ definition, and the oracle reformulation is the relativization. The lemma is the hinge of the following bridge.

\begin{thm}[Limiting recursion and learnability {\cite{Gold1965,Putnam1965,Shoenfield1959}}]
\label{thm:limit-bridge}
A family of r.e.\ sets $\mathcal{L} = \{L_i\}_{i \in I}$ is identifiable in the limit (Gold) if and only if the index function $i \mapsto e_i$, returning a correct hypothesis index for each $L_i$, is limiting recursive; the classes identifiable in the limit are exactly those whose index sets are $\Delta^0_2$.
\end{thm}

On the computability side, $\Delta^0_2$ is the level computable with a halting oracle, equivalently the sets decidable by a procedure allowed finitely many mind changes (\cref{thm:shoenfield}). On the learning side, identification in the limit is exactly the requirement that the conjecture sequence stabilize, which is exactly limiting recursion. The learning half of the bridge is verified: explanatory (Gold) identifiability is equivalent to the existence of a learner converging with finitely many conjecture changes.\formal{FLT_Proofs.Theorem.Gold.mind_change_characterization} The number of changes is itself a complexity measure, recorded as an ordinal that is finite exactly when the learner converges correctly.\formal{FLT_Proofs.Complexity.MindChange.MindChangeOrdinal}

The equivalence settles three points used in \Cref{ch:gold}.

\begin{enumerate}[leftmargin=*,itemsep=3pt]
\item \textbf{Upper bound on learnability.} No class whose index set lies above $\Delta^0_2$ in the arithmetical hierarchy is identifiable in the limit. This bounds Gold-style learning from above.

\item \textbf{Mind changes as approximation stages.} The number of $s$ with $\hat{f}(x,s) \neq \hat{f}(x,s+1)$ is the number of mind changes the learner makes before converging. The mind-change ordinals of \Cref{ch:ordinals} refine this into a transfinite hierarchy.

\item \textbf{Identification in the limit is canonical.} Gold's criterion is the natural first level above the computable: computability-in-the-limit, equivalently $\Delta^0_2$-computability. It is not one convergence requirement among many but the $\Delta^0_2$ one.
\end{enumerate}

\section{What This Chapter Established}
\label{sec:ch3-summary}

The chapter fixed the definitions---DFA, NFA, PDA, Turing machine, CFG, G\"odel numbering, partial computable function---needed for the identification-in-the-limit chapters, each typed against the verified computability substrate, and proved the load-bearing facts in full. Two structural ideas carry forward.

\begin{enumerate}[leftmargin=*,itemsep=3pt]
\item \textbf{The \rel{restricts}/\rel{extends\_grammar} distinction.} The NFA--DFA pair (conservative, same power, \cref{thm:nfa-dfa}) and the PDA--DFA pair (generative, strictly more power, \cref{prop:anbn}) are the first concrete instance of an edge-type distinction that recurs throughout the graph. Whether a generalization is conservative or generative determines whether existing theory transfers or must be rebuilt.

\item \textbf{The $\Delta^0_2$ bridge.} Identification in the limit is a precise computability concept: computation in the limit, equivalent to $\Delta^0_2$-computability (\cref{thm:shoenfield,thm:limit-bridge}). The equivalence anchors Gold's model in the arithmetical hierarchy and grounds the mind-change hierarchy of \Cref{ch:ordinals}.
\end{enumerate}

\section{Exercises and Open Problems}
\label{sec:ch3-open}

Each problem is anchored to a concept above and tagged: a \emph{known unknown} (\textsc{ku}) is an articulated open question; an \emph{unknown known} (\textsc{uk}) is a result present in the literature or the verified development but not yet absorbed into a clean general statement. Verified handles are named where they exist.

\begin{openprob}[Formalizing the limit-lemma bridge, \textsc{uk}]
\label{op:bridge-formal}
\Cref{prop:limit-delta} proves one direction. The computability substrate types both sides (\leanname{FLT_Proofs.Computation.LimitingRecursive}, \leanname{FLT_Proofs.Computation.Delta02Class}) and verifies the learning half (\leanname{FLT_Proofs.Theorem.Gold.mind_change_characterization}). Formalize the converse of Shoenfield's limit lemma and assemble the full bridge of \cref{thm:limit-bridge} ($\mathrm{Ex}$-identifiable $\Leftrightarrow$ $\Delta^0_2$ index set) inside the development.
\end{openprob}

\begin{openprob}[Optimal NFA determinization blow-up, \textsc{ku}]
\label{op:subset-tight}
\Cref{thm:nfa-dfa} gives a DFA with at most $2^n$ states. Exhibit, for each $n$, an $n$-state NFA whose minimal equivalent DFA has exactly $2^n$ states (Moore, Meyer--Fischer), and formalize the matching lower bound on top of \leanname{FLT_Proofs.Computation.NFA'} and \leanname{FLT_Proofs.Computation.DFA'}, so that the conservative-generalization cost is pinned exactly.
\end{openprob}

\begin{openprob}[Myhill--Nerode and the exact regularity criterion, \textsc{uk}]
\label{op:myhill-nerode}
The pumping lemma (\cref{lemma:pumping}) is necessary but not sufficient for regularity. Formalize the Myhill--Nerode theorem---$L$ is regular iff its right-congruence has finite index, which also yields the unique minimal DFA---on top of \leanname{FLT_Proofs.Computation.IsRegularLang}. This is the structural foundation that the $L^\ast$ learner of \Cref{ch:exact} exploits.
\end{openprob}

\begin{openprob}[Arithmetical complexity of the criterion hierarchy, \textsc{ku}]
\label{op:criterion-hierarchy}
\Cref{thm:limit-bridge} places explanatory identification at $\Delta^0_2$. Locate the other Gold criteria---behaviorally correct, vacillatory, anomalous, finite---in the arithmetical hierarchy, and determine which strictly exceed $\Delta^0_2$. The criteria are typed (\leanname{FLT_Proofs.Criterion.Gold.BCLearnable} and siblings) but their hierarchy positions are not yet pinned.
\end{openprob}

\begin{openprob}[A functorial account of restricts vs.\ extends, \textsc{uk}]
\label{op:restricts-functor}
\Cref{def:restricts,def:extends-grammar} distinguish conservative from generative generalization informally. Make it precise: define a category of computational models and a notion of generalization morphism such that \rel{restricts} morphisms provably transport characterization theorems and \rel{extends\_grammar} morphisms provably need not, and decide whether the classification of a given pair is itself decidable.
\end{openprob}

\begin{openprob}[Learnability of context-free languages, \textsc{ku}]
\label{op:cfl-learnable}
$\{a^n b^n\}$ (\cref{prop:anbn}) is context-free. Characterize which subclasses of the context-free languages are identifiable in the limit from text and from informant, extending the text/informant separation of \Cref{ch:data} up one level of the Chomsky hierarchy, and relate the answer to the $\Delta^0_2$ ceiling of \cref{thm:limit-bridge}.
\end{openprob}

\begin{openprob}[The coding theorem for the computability substrate, \textsc{uk}]
\label{op:coding-theorem}
The substrate types Kolmogorov complexity and algorithmic probability (\leanname{FLT_Proofs.Computation.KolmogorovComplexity}, \leanname{FLT_Proofs.Computation.AlgorithmicProbability}). Formalize the coding theorem $K(x) = -\log \mathrm{P}(x) + O(1)$ relating them, with the constant explicit in the reference machine---the same gap left open for model selection in \Cref{ch:objects} (Problem in \S\ref{sec:ch1-open}).
\end{openprob}

\begin{openprob}[The transfinite mind-change hierarchy, \textsc{uk}]
\label{op:mind-change-ordinals}
The mind-change count is typed as an ordinal (\leanname{FLT_Proofs.Complexity.MindChange.MindChangeOrdinal}), finite exactly when the learner converges correctly. Formalize the Freivalds--Smith transfinite mind-change hierarchy---identification with at most $\alpha$ mind changes for a countable ordinal $\alpha$---and prove the strictness of the resulting hierarchy, which \Cref{ch:ordinals} develops informally.
\end{openprob}

\begin{openprob}[Closure of the recursive languages under learning operations, \textsc{ku}]
\label{op:recursive-closure}
\Cref{prop:closure} and \cref{tab:closure} record language-class closure. Determine the closure of the \emph{identifiable-in-the-limit} classes under union, intersection, and complement: if $\mathcal{C}_1$ and $\mathcal{C}_2$ are each identifiable from text, when is $\mathcal{C}_1 \cup \mathcal{C}_2$, and does the answer mirror the regular/context-free split of \cref{tab:closure}?
\end{openprob}

\begin{openprob}[Effective separation witnesses, \textsc{uk}]
\label{op:effective-witness}
\Cref{prop:anbn} witnesses $\mathrm{REG} \subsetneq \mathrm{CFL}$ by a single language. Give a uniform, effective procedure that, from a description of any two adjacent Chomsky levels, outputs a witness language separating them together with its non-membership proof, turning the proper-containment chain of \cref{fig:chomsky} into one constructive theorem rather than four independent separations.
\end{openprob}
